\documentclass[aspectratio=169,10pt,spanish]{beamer}

\input{../../../_preambulo_beamer.tex}
\input{../../../_comandos_beamer.tex}

\selectlanguage{spanish}

\title{MA1001B - Modelación Estadística para la Toma de Decisiones}
\subtitle{Lección 03: Tendencia central y dispersión}
\author{}
\institute{Tecnol\'ogico de Monterrey}
\date{}

\begin{document}

\begin{frame}[plain]
  \vspace{1.2cm}
  {\Large\bfseries MA1001B - Modelación Estadística para la Toma de Decisiones\par}
  \vspace{0.7cm}
  {\large Lección 03: Tendencia central y dispersión\par}
  \vspace{0.7cm}
  {\color{TecRojo}\rule{\textwidth}{1.2pt}}
  \vspace{0.8cm}

  Facultad MA1001B\\[0.3cm]
  Periodo\\[0.3cm]
  Tecnol\'ogico de Monterrey
\end{frame}

\begin{frame}{La pregunta descriptiva}
  \begin{alertblock}{Pregunta guía}
    ¿Qué resúmenes describen un tiempo de procesamiento típico y la
    variabilidad entre pedidos?
  \end{alertblock}

  \vspace{0.5em}
  Al final de esta lección, usted puede:
  \begin{itemize}
    \item calcular e interpretar la media y la mediana;
    \item calcular la desviación estándar muestral y el rango intercuartílico;
    \item emparejar una medida de centro con una medida compatible de
      dispersión;
    \item explicar cómo un valor extremo afecta cada resumen.
  \end{itemize}

  \vfill
  \scriptsize Los 240 registros de pedidos son sintéticos. No se usan datos reales de empresa.
\end{frame}

\begin{frame}{Centro y dispersión}
  \small
  \begin{center}
    \begin{tabular}{p{0.18\textwidth}p{0.31\textwidth}p{0.37\textwidth}}
      \toprule
      \textbf{Resumen} & \textbf{Cálculo} & \textbf{Pregunta que responde} \\
      \midrule
      Media & $\bar{x}=\frac{1}{n}\sum x_i$ & ¿Dónde se equilibran los valores? \\
      Mediana & Valor central ordenado & ¿Qué valor parte los datos ordenados a la mitad? \\
      $s$ muestral & $\sqrt{\frac{\sum(x_i-\bar{x})^2}{n-1}}$ & ¿Cuánto varían los valores alrededor de la media? \\
      IQR & $Q_3-Q_1$ & ¿Qué tan ancho es el 50\% central? \\
      \bottomrule
    \end{tabular}
  \end{center}

  \vspace{0.8em}
  \begin{exampleblock}{Regla de interpretación}
    Reporte los resúmenes en contexto y en las unidades de la variable.
    Empareje media con $s$ muestral, o mediana con IQR, para que centro y
    dispersión cuenten una historia coherente.
  \end{exampleblock}
\end{frame}

\begin{frame}[fragile]{Paso 1: Media y mediana}
  \begin{columns}[T]
    \begin{column}{0.40\textwidth}
      \small
      \begin{lstlisting}
values = sample[
    "processing_time_minutes"
]

mean = values.mean()
median = values.median()
      \end{lstlisting}

      \begin{block}{Salida verificada}
        $n=240$ pedidos\\
        Media $=45.19$ min\\
        Mediana $=44.70$ min\\
        Diferencia $=0.49$ min
      \end{block}
    \end{column}
    \begin{column}{0.60\textwidth}
      \centering
      \includegraphics[width=\textwidth]{../figuras/estadistica_descriptiva_01_centro.png}
      \vspace{0.15em}
      \scriptsize Ambos números describen el centro. Ninguno describe la
      dispersión.
    \end{column}
  \end{columns}
\end{frame}

\begin{frame}{Paso 2: Dos medidas de dispersión}
  \begin{columns}[T]
    \begin{column}{0.39\textwidth}
      \small
      Desviación estándar muestral:
      \[
        s=\sqrt{\frac{\sum(x_i-\bar{x})^2}{n-1}}=6.38\text{ min}
      \]

      Intervalo de cuartiles:
      \[
        Q_1=40.53,\quad Q_3=49.01
      \]
      \[
        \mathrm{IQR}=Q_3-Q_1=8.48\text{ min}
      \]

      $s$ mide la variación alrededor de la media. El IQR cubre la mitad
      central de las observaciones ordenadas.
    \end{column}
    \begin{column}{0.61\textwidth}
      \centering
      \includegraphics[width=\textwidth]{../figuras/estadistica_descriptiva_02_dispersion.png}
    \end{column}
  \end{columns}
\end{frame}

\begin{frame}{Paso 3: Sensibilidad a un valor extremo}
  \small
  \textbf{El pedido 300193 cambia de 65.22 a 180.00 minutos.}

  \vspace{0.1em}
  \centering
  \includegraphics[width=0.80\textwidth]{../figuras/estadistica_descriptiva_03_sensibilidad.png}

  \vspace{0.1em}
  \begin{tabular}{lrrrr}
    \toprule
    \textbf{Resumen} & Media & Mediana & $s$ muestral & IQR \\
    \textbf{Cambio (min)} & $+0.48$ & $+0.00$ & $+4.34$ & $+0.00$ \\
    \bottomrule
  \end{tabular}

  \vspace{0.1em}
  \scriptsize Los cambios mostrados están redondeados a dos decimales. La
  sensibilidad por sí sola no diagnostica una anomalía.
\end{frame}

\begin{frame}{Elegir resúmenes descriptivos}
  \small
  \textbf{Media con desviación estándar muestral}\par\vspace{0.2em}
  Este par usa cada observación. Suele describir bien una distribución
  aproximadamente simétrica cuando ningún extremo domina los resúmenes.

  \vspace{0.7em}
  \textbf{Mediana con rango intercuartílico}\par\vspace{0.2em}
  Este par depende de posiciones ordenadas. Suele resistir mejor la
  influencia de un número pequeño de valores extremos.

  \vspace{0.8em}
  \begin{alertblock}{Una descripción defendible}
    Inspeccione la distribución, enuncie las unidades y explique la elección.
    Estos emparejamientos son pautas útiles, no reglas automáticas.
  \end{alertblock}

  Un identificador numérico como \texttt{order\_id} sigue etiquetando un
  registro. No es una medición cuantitativa para resumir con media o
  desviación estándar.
\end{frame}

\begin{frame}{Verificación de conceptos}
  \begin{enumerate}
    \item ¿Qué preguntas distintas responden la media y la desviación
      estándar muestral sobre el tiempo de procesamiento?
    \item Con $Q_1=40.53$ y $Q_3=49.01$, calcule e interprete el IQR.
    \item ¿Qué par centro-dispersión describe mejor datos afectados por unos
      pocos retrasos extremos, y por qué?
    \item ¿Por qué el valor de 180 minutos requiere más evidencia antes de
      etiquetarlo como anomalía o error de registro?
  \end{enumerate}

  \vspace{0.5em}
  \begin{block}{Conexión con la Lección 04}
    Los boxplots y las reglas explícitas de detección añaden una forma
    reproducible de marcar valores para investigación. El contexto de
    negocio sigue determinando qué significa la bandera.
  \end{block}

  \vfill
  \scriptsize Fuente: OpenStax, \textit{Introductory Business Statistics 2e},
  Capítulo 2, Secciones 2.2, 2.3 y 2.7. CC BY-NC-SA 4.0.
\end{frame}

\appendix



\end{document}
